% Preface: the README's About and "There Is No Instrument Here" sections, info/README.md's Who
% This Is For, Prerequisites, course plan and What This Edition Leaves Out, info/toolkit.md, and
% lectures/README.md's account of the exercises, as the front matter of a book.

\chapter{Preface}
\markboth{Preface}{Preface}

This book is for embedded and software engineers who read a schematic every week and have never
been shown how one is arrived at. You can write C++, you are comfortable rearranging an
expression, and you have at some point been handed an amplifier stage with no way to tell whether
it was right. It takes you from Ohm's law to a discrete operational amplifier you have designed,
sized, predicted and solved, transistor by transistor.

\textbf{The subject is the transistor amplifier.} Circuit theory, passives and the operational
amplifier are built in the first four chapters because the last six need them, not for their own
sake. That makes the book deliberately narrow: ten hours of lectures is not long enough for
everything an analog textbook carries, so this edition keeps what decides a design and says
plainly what it dropped, in \prefref{What this edition leaves out}{pre:sec:cuts}.

\section*{How the book is organised}
The book is typeset from a ten-lecture course, and each chapter is one lecture. Its sections are
the lecture's appendices, in the same order and with the same subsections, so a chapter here and a
lecture directory in the companion repository can be read side by side:

\begin{booktable}{@{}p{4.6em}p{3.4em}p{5.6em}L@{}}
  \thead{Chapter} & \thead{Lecture} & \thead{Part} & \thead{Topic} \\
  \midrule
  1  & \file{L01} & Foundations & Circuits, units, and the nodal solver \\
  2  & \file{L02} & Foundations & Reactance, phasors, and frequency response \\
  3  & \file{L03} & Foundations & Passive filters and the operational amplifier \\
  4  & \file{L04} & Foundations & Feedback, active filters, and the diode \\
  5  & \file{L05} & Transistors & The transistor as a device and as a switch \\
  6  & \file{L06} & Transistors & Biasing, and what an emitter resistor actually buys \\
  7  & \file{L07} & Transistors & Small-signal analysis: $r_e$, the emitter factor, the cascode \\
  8  & \file{L08} & Transistors & Followers, and the output stage they turn into \\
  9  & \file{L09} & Transistors & The differential amplifier \\
  10 & \file{L10} & Transistors & Building an operational amplifier, and the capstone \\
\end{booktable}

\noindent Every chapter opens with what is in it, works through the material, closes with a review
of what you should now be able to explain, and ends with its exercises. The appendices hold the
two written papers and their model answers.

\textbf{Three things about the ordering are deliberate}, and all three differ from how the subject
is usually taught:
\begin{itemize}
  \item \textbf{The solver comes first.} Chapter~\ref{c1:ch:circuits} has you write nodal analysis
    before you have met a single active device. It costs an hour and it pays for the rest of the
    book: from then on every claim the material makes about a circuit can be checked against a
    program you wrote, and no result has to be taken on authority.
  \item \textbf{The operational amplifier comes before the transistor.} You use one as a black box
    in Chapters~\ref{c3:ch:filters} and \ref{c4:ch:feedback}, and only in
    Chapter~\ref{c10:ch:amplifier} do you build one. That inversion is the whole arc of the book:
    the thing you were told to trust in the fourth chapter is the thing you design in the last.
  \item \textbf{The imperfections come before the topologies.} Biasing, thermal drift and the
    emitter factor are met in Chapters~\ref{c6:ch:bias} and \ref{c7:ch:smallsignal}, before any of
    the clever stages. Meeting them early means that the question ``why would anyone use a current
    mirror as a load'' has an answer you can compute rather than remember.
\end{itemize}

\section*{What you should be able to do at the end}
Having worked through the book, you should be able to:
\begin{itemize}
  \item Solve any resistive network by nodal analysis, by hand and in code.
  \item Read and predict a first-order frequency response, magnitude and phase, from its corner
    alone, and say what a source resistance does to it.
  \item Apply the two operational amplifier rules to any standard configuration, and say precisely
    where each of them fails.
  \item Compute a loop gain and everything it decides: gain error, distortion, impedances,
    bandwidth.
  \item Bias a transistor stage to a specification, including the base current's effect on the
    divider that sets it, and state its thermal drift.
  \item Derive gain, input resistance and output resistance for any stage in the book from one
    small-signal model, and say which node each result belongs to.
  \item Choose between a follower, a Darlington and a redesign from the arithmetic rather than
    from habit.
  \item Size a class-AB output stage and say why its bias generator is bolted to a heatsink.
  \item Design a differential pair to a rejection requirement, and say which components do not
    affect the answer.
  \item Compute an operational amplifier's gain budget with loading, close the loop around it, and
    predict the error.
\end{itemize}

\phantomsection
\section*{What you are assumed to know}
\label{pre:sec:prerequisites}
You are expected to arrive already comfortable with:
\begin{itemize}
  \item \textbf{Ohm's law}, and what a resistor, a capacitor and a voltage source are. Nothing
    beyond that is assumed about circuits; Chapters~\ref{c1:ch:circuits} and \ref{c2:ch:reactance}
    build the rest.
  \item \textbf{Engineering calculations}: rearranging an expression, working in decibels, taking
    a logarithm, and being willing to carry an approximation and say how wrong it is.
  \item \textbf{Complex numbers} far enough to accept that $j$ rotates by ninety degrees.
    Chapter~\ref{c2:ch:reactance} uses that and builds the frequency domain out of it.
\end{itemize}

\textbf{No semiconductor physics is taught, and none is needed.} The transistor is introduced as a
device with an exponential, a current gain and a saturation voltage. Depletion regions, carrier
transport and the Ebers-Moll derivation are named once and not used.

\textbf{C++ is used, and this is not a C++ book.} You write a small analysis library across the
ten chapters. The C++ involved is free functions, plain structs, \code{std::vector} and
\code{std::complex}. There is no class hierarchy anywhere in it. If you have written embedded
C++, you have written harder code than any of it.

\phantomsection
\section*{There is no instrument here, and that is the design}
\label{pre:sec:instrument}
A book like this one usually comes with a breadboard, a bench supply, an oscilloscope and LTspice.
This one does not. There is nothing to install but a C++ compiler, \code{make} and \code{git}.

What replaces it is that \textbf{you write the instrument}. Across the ten chapters you build
\code{ael}, an analog analysis library in C++: a netlist, a nodal solver, complex stamps for
reactance, a Newton-Raphson loop, device models for the diode, the BJT and the MOSFET, and on top
of them the closed-form results for every amplifier stage in the book. By Chapter~\ref{c10:ch:amplifier}
they compose into one program that takes a discrete operational amplifier, predicts every node in
it, and then solves it.

\textbf{None of that code is in this book or in the repository.} Every component is specified in
prose, with its interface declaration, in the section of the chapter that asks for it, and every
one ships with a test suite that judges what you wrote.

That trade cuts both ways:
\begin{itemize}
  \item \textbf{What you lose.} No breadboard, no oscilloscope, no real device spread, no
    parasitics, no thermal behaviour you can feel with a fingertip, and no experience driving an
    instrument. Where a number here differs from what a real measurement would give, the text says
    so and says by how much.
  \item \textbf{What you gain.} Every step is a function you wrote and can read. ``The simulator
    says the gain is $-84$'' stops being an oracle. You know which term that came from, what it
    assumed, and what would have to change to move it.
\end{itemize}

Verification works in three layers, and none of them is a bench:
\begin{itemize}
  \item \textbf{The Cross-check.} Every chapter ends with one exercise where you compute a number
    by hand, then run your own closed-form code on it, then solve the same circuit numerically
    with your own nodal solver. Three legs, and they are not equally authoritative: the hand
    calculation drops terms on purpose, the closed form drops fewer, and the solver drops none but
    tells you nothing about why. When two of them disagree by 8 per cent, finding out which one is
    wrong, or that neither is, is the exercise, and each one ends by saying how large a
    discrepancy to expect.
  \item \textbf{The shipped test suites.} Every component you implement comes with a unit-test
    suite written against the QAcademy Test framework. They tell you \emph{whether} you are right
    without telling you \emph{why}, which is the middle ground self-study otherwise lacks.
  \item \textbf{The two models against each other.} From Chapter~\ref{c7:ch:smallsignal} on, every
    stage is computed twice: once from the closed form, once by the solver from the same netlist.
    The closed form is the one you can reason with and the solver is the one that is right. Where
    they part company, the text names the term responsible.
\end{itemize}

\phantomsection
\section*{The toolkit you build}
\label{pre:sec:toolkit}
\code{ael} lives wherever you like, as long as you tell the test suites:

\begin{shell}
export AEL_DIR=~/ael
make test
\end{shell}

\noindent The layout every suite assumes is \file{$AEL_DIR/include/ael/...} for the headers and
\file{$AEL_DIR/source/ael/...} for the translation units. \textbf{The header path is part of each
specification, not a suggestion}: the suites include exactly these paths, and a component written
somewhere else stays dormant however correct it is.

\begin{booktable}{@{}p{4.4em}p{12.4em}L@{}}
  \thead{Chapter} & \thead{Header} & \thead{What it does} \\
  \midrule
  1  & \code{ael/net/netlist.hpp}      & \code{Netlist}: nodes, resistors, and independent sources \\
  1  & \code{ael/mna/solver.hpp}       & \code{solve}: the DC nodal solve \\
  2  & \code{ael/ac/sweep.hpp}         & \code{solveAt} and \code{sweep}: the same solve over complex admittances \\
  3  & \code{ael/filter/response.hpp}  & Filter responses, resonance, Q, and the loaded cascade \\
  3  & \code{ael/opamp/ideal.hpp}      & The configuration gains and the Schmitt thresholds \\
  4  & \code{ael/feedback/loop.hpp}    & Loop gain, closed-loop gain, and the error \\
  4  & \code{ael/device/diode.hpp}     & The diode equation, its conductance, and the step limiter \\
  4  & \code{ael/nr/solve.hpp}         & The nonlinear DC operating point: the linear solve, in a loop \\
  5  & \code{ael/device/bjt.hpp}       & The transport model, the three regions, and the switch \\
  5  & \code{ael/device/mosfet.hpp}    & The square law, its two regions, and $g_m$ \\
  6  & \code{ael/bias/point.hpp}       & The quiescent point including the droop, and the drift \\
  7  & \code{ael/ssm/model.hpp}        & The $r_e$ model: gain, both resistances, EF, the cascode \\
  8  & \code{ael/follower/stage.hpp}   & Follower gain, both resistances, and the Darlington \\
  8  & \code{ael/output/classab.hpp}   & The 26 mV rule, the bias voltage, and the idle current \\
  9  & \code{ael/diffpair/pair.hpp}    & Both gains, the rejection ratio, the tanh, and the offset \\
  10 & \code{ael/report/amplifier.hpp} & The capstone: the budget, the open loop, the closed loop \\
\end{booktable}

\noindent Namespaces mirror the header \emph{directory}, so everything in \code{ael/ssm/model.hpp}
is in \code{ael::ssm}. Two directories hold more than one component and take the file name as a
third level: \code{ael/device/bjt.hpp} is \code{ael::device::bjt}, and \code{ael/opamp/ideal.hpp}
is \code{ael::opamp::ideal}. Units are SI and unqualified throughout: where the text quotes 26 mV,
the code carries \code{0.026}.

\begin{aside}
\textbf{Two sign conventions, fixed once}, because every component in the book reads them. A
current source \emph{injects at its second node}: \code{addCurrentSource(from, to, current)} takes
current out of \code{from} and pushes it into \code{to}, so one milliamp into a node with one
kilohm to ground puts that node at $+1$ V. And a voltage source \emph{reports the current leaving
its positive terminal}, so a 10 V source driving one kilohm reports $+10$ mA --- which is the
opposite sign to the raw modified-nodal unknown. Get either backwards and \secref{c1:sec:b5}'s
suite says so on the first run; get either backwards and keep it, and every bias calculation in
Part~2 comes out inverted.
\end{aside}

\section*{What the exercises look like}
Six to ten per chapter, each labelled with its kind:

\begin{booktable}{@{}p{9.6em}L@{}}
  \thead{Kind} & \thead{What it asks of you} \\
  \midrule
  \kindtag{Recall}                     & State a result without looking it up. \\
  \kindtag[fillaccent2]{Hand calculation} & Produce a number with a pencil, making the
    approximations deliberately and saying what each one costs. \\
  \kindtag[fillaccent4]{Design}        & Start from a requirement and end with E12 component
    values, the error they leave, and whether you would accept it. \\
  \kindtag[fillaccent]{Code}           & Write a component of the toolkit, specified in prose and
    judged by the shipped suite. \\
  \kindtag[fillaccent3]{Cross-check}   & Compute one number three ways and reconcile them. Exactly
    one per chapter, and it is always the last. \\
\end{booktable}

\textbf{The Cross-check is the signature exercise.} You compute a number by hand, run your own
closed-form code on the same problem, and then solve the same circuit numerically with your own
nodal solver. Three legs, none of them authoritative on its own. Every Cross-check ends by saying
how large a discrepancy to expect and what each size of discrepancy would mean, because ``your two
answers differ by 8 per cent'' is the lesson rather than a failure.

\textbf{The solutions are not in this book.} They are published in full in the companion
repository, under \file{lectures/LNN/appendix}, including the plausible wrong answer where there is
one --- this is a self-study course, and a reader with nobody to ask cannot be left with an
unanswered exercise. Try each exercise before you read its solution.

\section*{Two written papers, and what they are not}
Nothing in the course this book comes from is marked. Assessment is the exercises in every chapter,
the test suite each component comes with, and the Cross-check that closes every chapter.

Appendices~\ref{exam:ch:about} to \ref{ansb:ch:answers} hold two four-hour papers with worked
solutions, and they check something else: \textbf{what you can reconstruct on paper, with the book
out of the room.} Ten questions each, one per chapter.

\textbf{They exist purely so that you can test your own knowledge after working through the book.
They gate nothing, they are not a qualification, and no part of the book requires them.}

\textbf{Take one after the last chapter.} Every paper has one question per chapter, so sitting one
partway through examines material you have not read yet.

\phantomsection
\section*{What this edition leaves out}
\label{pre:sec:cuts}
About a thousand pages of standard treatment became ten chapters. What went, in full, and where it
is:

\begin{booktable}{@{}Lp{14.4em}@{}}
  \thead{Cut} & \thead{Kept here as} \\
  \midrule
  Basic electronics as eight chapters: transformer design, rectifier sizing, rms &
    Chapters~1 and 2, at the level later work needs \\
  Passive filters as six chapters, each with its own full impedance treatment &
    Chapter~3, one treatment covering all of them \\
  Telescopic cascodes and CMOS amplifier topologies &
    Not covered; the cascode itself is \secref{c7:sec:b5} \\
  Output stages up to 1 kW: triple class-AB, CFP-EF, Zobel networks, protection &
    \Secref{c8:app:c}, as reading \\
  Sixteen appendices of full derivations &
    The results, with the derivations cited \\
  Noise analysis, Miller's theorem as a chapter, and CMOS frequency response &
    The Miller effect in \secref{c7:sec:b5}; noise not at all \\
  Oscillators, voltage regulators, and a CMOS comparator &
    Not covered \\
\end{booktable}

\noindent The last two rows are material a full treatment would carry and this edition does not.
\textbf{Noise in particular is a real omission:} it decides the input stage of any amplifier that
matters, and it is the first thing to add if this course is ever extended. PCB layout, datasheet
reading and phase margin as a design procedure are in neither.

\textbf{Where the obvious answer is wrong}, the correction is taught rather than applied quietly.
The two largest are Chapter~\ref{c6:ch:bias}'s thermal argument and
Chapter~\ref{c7:ch:smallsignal}'s output resistance, and both are worked demolitions: the tempting
claim, the arithmetic that breaks it, and the version that survives.

\section*{Further reading}
There is no set textbook: every derivation, figure and number here is written for this course. Two
books are worth having beside it.
\begin{itemize}
  \item \textit{The Art of Electronics}, Horowitz and Hill, 3rd edition, is the natural companion.
    It uses the same intrinsic-emitter-resistance approach this book is built on, and each chapter
    names the parts of it that cover the same material.
  \item \textit{Microelectronic Circuits}, Sedra and Smith, if you want the hybrid-pi treatment
    instead, and the device physics this book skips.
\end{itemize}

\section*{Setting up}
Everything here runs in a Linux terminal: on Linux itself, or on Windows through WSL. You need a
C++17 compiler, \code{make} and \code{git}, and nothing else --- no LTspice, no MATLAB, no
breadboard, no instrument. On WSL or Ubuntu:

\begin{shell}
sudo apt -y update
sudo apt -y install git make g++
g++ --version             # 9 or newer; CI builds with the g++ on ubuntu-24.04.
\end{shell}

\noindent Then clone the companion repository, together with the test framework it keeps as a
submodule:

\begin{shell}
git clone --recursive https://github.com/qrtech-academy/analog-electronics.git
\end{shell}

\noindent Three commands at its root cover everything the book asks of it:

\begin{shell}
make help                 # List every target.
make test                 # Build and run each chapter's suite against your toolkit.
make lint                 # Links, house rules, quoted numbers, and formatting.
\end{shell}

\noindent \code{make test} needs \code{AEL_DIR} to know where your toolkit is. \textbf{With it
unset the suites still build and still run}, and report that only each chapter's reference tests
were active: every test for a component you have not written yet sits behind an
\code{#if __has_include}, so it switches itself on the moment the header appears. That is the
correct state on day one, and \textbf{anything not yet implemented is skipped loudly rather than
failing}, which is what makes \code{make test} meaningful on the first day and on the last.
